\documentclass[11pt]{article}
\usepackage[margin=.77in]{geometry}
\usepackage{amsmath,amssymb,amsthm,booktabs,hyperref}
\hypersetup{colorlinks=true,urlcolor=blue,linkcolor=black,citecolor=black}
\newtheorem{theorem}{Theorem}
\newtheorem{lemma}[theorem]{Lemma}
\newtheorem{corollary}[theorem]{Corollary}
\newcommand{\F}{\mathbb F}
\title{When cyclic labels $a_i^2+a_{i+1}$ cover every residue}
\author{Blake Stone's mathematics project\\\small Proof developed with OpenAI Codex}
\date{November 26, 2026}
\begin{document}\maketitle
\begin{abstract}
There is a circular ordering of $\mathbb Z/n\mathbb Z$ in which the labels $a_i^2+a_{i+1}$ also run through $\mathbb Z/n\mathbb Z$ exactly once if and only if $\gcd(n,6)=1$. This proves Zhi-Wei Sun's 2013 conjecture for all primes $p>3$ and describes the support of the counting sequence OEIS A229827 for every modulus. The prime-field argument constructs two permutation functions with the required difference, translates one of them, and joins its cycles by switches at opposite arguments. Polynomial identities in the translation parameter guarantee that the switches can connect everything. A digit lift and the Chinese remainder theorem handle arbitrary moduli.
\end{abstract}

Zhi-Wei Sun asked whether, for every prime $p>3$, there is a single-cycle permutation $\sigma$ of $\F_p$ such that $x\mapsto x^2+\sigma(x)$ is also a permutation \cite{oeis}. His sequence counts such cycles on all of $\mathbb Z/n\mathbb Z$. The convention of listing residues as $1,\ldots,n$ or $0,\ldots,n-1$ gives the same count. His broader field and subset conjectures \cite{slides} are considered separately in the companion note.

\begin{theorem}\label{th:main}
For every $n\ge2$, a circular ordering $(a_i)$ of all residues modulo $n$ with all $n$ labels $a_i^2+a_{i+1}$ distinct exists exactly when $\gcd(n,6)=1$. In particular $\mathrm{A229827}(n)>0$ exactly on these indices.
\end{theorem}
If the two lists exhaust the same group, subtracting the sums of their entries gives $\sum_{x=0}^{n-1}x^2\equiv0\pmod n$. The square-sum identity makes this equivalent to $6\mid(n-1)(2n-1)$, or $\gcd(n,6)=1$. The remaining sections prove sufficiency.

\section{Joining cycles without changing the set of labels}
Let $X$ be a finite set equipped with an involution $J$. Let $f$ be a permutation such that its labels $P(x)+f(x)$ are distinct in an additive group, where $X$ is a subset of that group and $P$ is constant on each $J$-orbit.
\begin{lemma}\label{lem:splice}
If the graph on $X$ whose edges join $x$ to $f(x)$ and to $J(x)$ is connected, swaps of successor values at pairs $x,J(x)$ produce a single-cycle permutation with the same label multiset.
\end{lemma}
\begin{proof}
Contract the cycles of $f$ in the indicated undirected graph. Choose a spanning tree of the connected quotient. For every edge of the tree, exchange the successors at its two marked arguments $x,J(x)$. When that edge is used, the two current cycles are distinct, since earlier edges of a tree cannot have connected its endpoints. Exchanging successors between two distinct cycles combines them into one cycle. The two labels exchange places because $P(x)=P(J(x))$. All other labels are untouched. At the end one cycle remains.
\end{proof}

\section{The prime field}
Fix a prime $p>3$, put $m=(p-1)/2$, and write $\chi$ for the quadratic character with $\chi(0)=0$. There are nonzero $k,\ell$ with $k+\ell=-1$ and $\chi(k)=-\chi(\ell)$. Indeed $\sum_z\chi(z(z+1))=-1$. To see this, count solutions of $w^2=z(z+1)$: the change of variables $(2z+1-2w)(2z+1+2w)=1$ puts them in bijection with $\F_p^\times$. Thus among $z\ne0,-1$ there are $(p-2+\chi(-1))/2>0$ with $\chi(-z(z+1))=-1$. Take any such $k=z$, $\ell=-1-z$.

Choose one representative $x$ from each nonzero pair $\{x,-x\}$. Set $f(0)=0$ and
\[
 f(x)=kx^2,\qquad f(-x)=\ell x^2\quad\text{for each representative }x.
\]
The squares run once through the nonzero squares as the representatives vary. Their two multiples by $k,\ell$ divide $\F_p^\times$ into its two square classes, so $f$ is bijective. The labels $x^2+f(x)$ and $x^2+f(-x)$ have multipliers $1+k=-\ell$ and $1+\ell=-k$; they also give a permutation. Hence for every $t\in\F_p$ both $f_t=f+t$ and $x\mapsto x^2+f_t(x)$ are permutations.

\begin{lemma}\label{lem:moment}
Let $d$ be the least positive integer with $k^d+\ell^d=0$ in $\F_p$. For at least $p-d$ elements $t\in\F_p$, the graph generated by $x\mapsto f_t(x)$ and $x\mapsto-x$ is connected. In particular at least $(p+1)/2$ translations work.
\end{lemma}
\begin{proof}
Euler's criterion and the opposite characters show $k^m+\ell^m=0$, so $d\le m$. If a graph is disconnected, take one component $B$, a nonempty proper subset of $\F_p$. Both $f_t$ and negation preserve $B$. Write $M_j=\sum_{x\in B}x^j$, using $M_0=|B|\ne0$ in the prime field. Odd moments vanish by negation. For $r\ge1$, invariance under $f_t$ gives
\begin{equation}\label{mom}
 \frac{k^r+\ell^r}{2}M_{2r}
 =\sum_{x\in B} f(x)^r
 =\sum_{x\in B}(x-t)^r
 =\sum_{j=0}^{\lfloor r/2\rfloor}\binom r{2j}(-t)^{r-2j}M_{2j}.
\end{equation}
The first equality uses pairs $x,-x$; the possible zero element contributes zero to both sides.

Put $A_0(T)=1$. For $1\le r<d$, define $A_r(T)$ as the right side of \eqref{mom} with $t,M_{2j}$ replaced by $T,A_j(T)$, divided by $(k^r+\ell^r)/2$. This is recursive since $\lfloor r/2\rfloor<r$. Induction gives $M_{2r}=M_0 A_r(t)$ and $\deg A_r\le r$. At $r=d$ the left side of \eqref{mom} is zero. Dividing by $M_0$ gives $D(t)=0$ for the fixed polynomial
\[
 D(T)=\sum_{j=0}^{\lfloor d/2\rfloor}\binom d{2j}(-T)^{d-2j}A_j(T).
\]
Its first term is $(-T)^d$; every other term has degree at most $d-j<d$. Thus $D$ has at most $d$ roots. Every disconnected translation had to be one of these roots, proving the assertion.
\end{proof}

Take any connected translation supplied by Lemma~\ref{lem:moment} and apply Lemma~\ref{lem:splice} with $P(x)=x^2$ and $J(x)=-x$. This proves Theorem~\ref{th:main} for every prime $p>3$. The use of a prime field in $M_0\ne0$ is essential to this argument. The quadratic assertion on $\F_8$ fails as shown in the companion paper.

\begin{corollary}\label{cor:count}
For every prime $p>3$ the number of directed circular orders, modulo rotation, satisfies
\[
    \mathrm{A229827}(p)\ \ge\ \frac{(p^2-1)(p-3)}{16}.
\]
\end{corollary}
\begin{proof}
The same construction and moment proof work for any unordered pair
$\{k,\ell\}\subset\F_p\setminus\{0,-1\}$ such that $k,\ell$ have opposite quadratic characters and $k+1,\ell+1$ also have opposite characters. This is precisely what makes both initial tables permutations, and the moment argument uses only the first opposition. Each pair gives at least $(p+1)/2$ orders. From the resulting successor one recovers $t=\sigma(0)$, since zero is never switched, and then $\{k,\ell\}=\{\sigma(1)-t,\sigma(-1)-t\}$. Distinct parameter pairs and translations therefore give distinct directed circular orders.

To count the unordered pairs, write $C=p-2$ and $\epsilon=\chi(-1)$. Among $z\ne0,-1$ the number with $(\chi(z),\chi(z+1))=(a,b)$ is
\[
     N_{ab}=\tfrac14(C-a\epsilon-b-ab),\qquad a,b\in\{1,-1\}.
\]
This follows by expanding $\sum_{z\ne0,-1}(1+a\chi(z))(1+b\chi(z+1))/4$ and using the character sum above. The number of opposite-pattern pairs is
$N_{++}N_{--}+N_{+-}N_{-+}=(C^2-1)/8=(p-1)(p-3)/8$. Multiplication by $(p+1)/2$ proves the bound.
\end{proof}

\section{Lifting to powers and products}
Suppose $M=p^e$, $p>3$, and $\sigma$ is a full cycle with the desired labels modulo $M$. Write its orbit using least representatives $a_0,\ldots,a_{M-1}$ and read indices cyclically. Each lift to $\mathbb Z/pM\mathbb Z$ has a unique coordinate pair $(a_i,y)$ representing $a_i+My$, $y\in\F_p$. Define a successor by
\[
  (a_i,y)\longmapsto (a_{i+1},c_i y+b_i).
\]
Use $c_i=-1$ for all $a_i\equiv-1/2\pmod p$ and also for $a_i=0$, and $c_i=1$ otherwise. The first class occurs $p^{e-1}$ times. Thus $\prod_i c_i=1$. Choose $b_{M-1}=1$ and all other $b_i=0$; after one lap beginning at phase zero, $y$ becomes $y+1$. The full orbit therefore visits every lift exactly once.

Put $u_i$ equal to the least residue of $a_i^2+a_{i+1}$ modulo $M$ and write $a_i^2+a_{i+1}=u_i+M v_i$. The lifted labels in phase $i$ have first coordinate $u_i$ and top digit
\[
 v_i+b_i+(2a_i+c_i)y\pmod p.
\]
Here $M^2$ vanishes modulo $pM$. The coefficient is nonzero: it is $-2$ at $a_i\equiv-1/2$, $-1$ at $a_i=0$, and $2a_i+1\ne0$ elsewhere. Since each $u_i$ occurs at one phase, the lifted labels exhaust all residues. This proves the prime-power claim by induction.

Finally if $u,v$ are coprime and each admits a good cycle, take the product of the two successor permutations under the Chinese remainder theorem. Every orbit has length $\operatorname{lcm}(u,v)=uv$, and the labels in the two components are independent permutations. This proves Theorem~\ref{th:main}.

The construction also gives, for each $M=p^e$, $p>3$,
\[
    \mathrm{A229827}(pM)\ \ge\ (p-1)p^{M-1}\,\mathrm{A229827}(M).
\]
Indeed, fix $a_0=0$. For any $b_0,\ldots,b_{M-2}$, the last offset $b_{M-1}$ can take $p-1$ values for which the net displacement after one lap is nonzero. Any such displacement generates the additive group of $\F_p$, and the coefficients of the edge labels are unaffected. Different offsets change a successor at some $(a_i,0)$, and the projected successors distinguish different base circles. The same recovery of the coordinate cycles gives $\mathrm{A229827}(uv)\ge\mathrm{A229827}(u)\mathrm{A229827}(v)$ for coprime $u,v\ge2$.

\begin{corollary}
For $n$ coprime to six and any polynomial $Q(z)=a z^2+bz+c$ over $\mathbb Z/n\mathbb Z$ with $a$ a unit, some circular ordering has all labels $Q(z_i)+z_{i+1}$ distinct. Over each prime field of size $p>3$ this holds for every polynomial of degree at most two other than $c-X$.
\end{corollary}
\begin{proof}
Put $r=a^{-1}$ and $s=-b/(2a)$. Then $Q(rx+s)=r x^2+C$ for a constant $C$. Transport an order $(x_i)$ of Theorem~\ref{th:main} by $x_i\mapsto rx_i+s$; its new labels are $r(x_i^2+x_{i+1})+C+s$. For the linear prime-field case take $\sigma(x)=x+1$. For $Q=c-X$ any full cycle has $Q(x)+\sigma(x)\ne c$.
\end{proof}

\begin{thebibliography}{9}
\bibitem{oeis} Z.-W. Sun, OEIS A229827, circular permutations with all quadratic edge labels distinct, September 2013. \url{https://oeis.org/A229827}.
\bibitem{slides} Z.-W. Sun, \emph{Determinants, Permutations and Additive Combinatorics}, lecture slides, conjecture dated 30 September 2013. \url{https://www.mathnet.ru/PresentFiles/21090/sun.pdf}.
\end{thebibliography}
\end{document}
